\chapter{Logic --- What Follows, What Is Declined}\label{ch:logic}

% =============================================================================
% Ch.2 Logic. [SYNTHESIS]. Re-states canonical claims at canonical tags.
% Primary sources: SPEC_FTD_FRAMEWORK_V1.md §2 (2.1-2.6);
%                  THEOREM_COMMUTATIVITY_INDEPENDENCE.md (FTD-0243);
%                  LEDGER rows FTD-0243/0255/0256/0227/0225/0226/0228/0253/0252.
% Reproduced derivation: the commutativity-independence argument, making the
%   distinction {q,p} != 0 yet [q,p] = 0 intuitive.
% =============================================================================

% ---- chapter mini-map (the Logic band of the Master DAG) -------------------
\begin{figure}[t]
  \centering
  \resizebox{0.8\linewidth}{!}{%
  \begin{tikzpicture}[node distance=0pt]
    \node[bandlbl] at (-1.4,0.0) {\textsc{logic}\\[-1pt]\tiny what follows};
    \node[thm]  (alg)  at (2.0,0.0)  {commutative\\algebra $A_5$};
    \node[seed] (fc1)  at (5.6,0.0)  {FC-1\\\tiny declines $M$};
    \node[seed] (fc2)  at (9.2,0.0)  {FC-2\\\tiny native arrow};
    \node[wall] (shrp) at (12.8,0.0) {sharpness\\\tiny open};
    % within-band reading order
    \draw[con] (alg) -- (fc1);
    \draw[wk]  (fc1) -- (fc2);
    \draw[wk]  (fc2) -- (shrp);
  \end{tikzpicture}}
  \caption[This chapter in the chain]{This chapter in the chain. The Logic band:
    the substrate exposes a strictly \emph{commutative} observable algebra $A_5$
    (theorem-grade), on the strength of which two framework commitments are
    \emph{declared}---FC-1 declines the quantum measurement import, FC-2 takes the
    arrow of time as native---and one kernel stays open (the \emph{sharpness} of
    the observer's blind spot). Shapes and rules carry epistemic status exactly
    as in Figure~\ref{fig:masterdag}.}
  \label{fig:minimap-logic}
\end{figure}

% ---- intuition primer ------------------------------------------------------
\section*{The move this layer makes: a classical world that declines the quantum import}
\addcontentsline{toc}{section}{The move this layer makes}

\noindent
Ontology said \emph{what exists}: a finite ternary lattice carrying two
orthogonal fields, with a forced causal cone. Logic asks the next question:
given those things, \emph{what follows}, and---just as importantly---\emph{what
does the framework decline to assume}? This is the layer where FTD takes its
sharpest, most consequential positions, and it takes them in the open.

The move is this. The substrate, run forward under its five postulates, exposes
a world whose observables behave \emph{classically}: every quantity you can read
off a configuration is an ordinary real-valued function, and ordinary functions
multiply commutatively. There is no hidden non-commutativity waiting to be
distilled by a clever limit. Quantum mechanics' defining feature---observables
that fail to commute, so that measuring one disturbs another---is therefore not
something the substrate \emph{contradicts}; it is something the substrate simply
does not contain. FTD's response is not to manufacture it and not to deny it,
but to \emph{decline} it: to commit, as a matter of framework definition, that
the commutative substrate is the whole of its ontology of observables.

That single move has a delicate companion. The substrate's dynamics are far
from trivial---the leapfrog time-step gives its phase space a genuine symplectic
structure, the same structure that in classical mechanics generates Hamiltonian
flow. So one is tempted to say ``but surely a non-trivial symplectic structure
\emph{is} quantum-like.'' The whole of this chapter's reproduced derivation
exists to defuse that temptation precisely. There are \emph{two} brackets at
play, and they are different objects: the Poisson bracket lives on the
\emph{functions} and is non-zero; the observable commutator lives on the
\emph{algebra's product} and is zero. The substrate is symplectically rich and
algebraically commutative at the same time, with no contradiction. The quantum
deformation that would turn the first bracket into the second requires an
external parameter the substrate does not supply.

A second decision rides on the same logical base. P5 makes the update a
\emph{function}, not an invertible map---determinism is not reversibility---and
the genuinely finite part of the substrate, the ternary state field, is provably
many-to-one. So the framework declines global time-reversal symmetry as well,
taking the arrow of time as native rather than emergent. Two declinations, one
classical-logical foundation. Everything below states each at exactly the
epistemic grade the ledgers record: the algebra result is a theorem; the two
declinations are \emph{declarations}, not derivations; and one piece of the
observer story is honestly left open.

% ---- tagged canonical content ----------------------------------------------
\section{Determinism is not reversibility}\label{sec:det-not-rev}

Postulate~P5 states that the update map is a deterministic function of the
configuration. It does \emph{not} state that the map is invertible. A
deterministic map may be many-to-one---heat diffusion is the canonical example
of a process that is deterministic \emph{and} irreversible. This distinction is
not a technicality; it carries the entire weight of FC-2 below
(\S\ref{sec:fc2}). \etag{syn}~(\ftdid{0256}, licensing boundary \ftdid{0253}).

The point to fix is that the postulate base \emph{permits} an arrow without yet
\emph{committing} to one. Whether reversibility is a property of the world is,
at the level of P1--P5, an open fork; \S\ref{sec:fc2} records which branch FTD
takes and why that is a declaration rather than a theorem.

\section{The substrate's observable algebra is commutative}\label{sec:commutative}

Let the \emph{observable algebra} $A_5$ be the smallest collection of
real-valued functionals of the configuration $(s, J, \mathrm{wave\_vel},
\mathcal{L})$ that contains those generators and is closed under pointwise
real-linear combination, pointwise product, Moore-neighbourhood sums, and
composition with the deterministic update map $U$. The central logical fact of
the substrate is:

\begin{result}[Commutativity of $A_5$]\label{res:commutative}
The observable algebra $A_5$ is strictly commutative.
\end{result}

\noindent
\etag{thm}~(\ftdid{0243}; algebraic core machine-checked in
\texttt{lean/Standalone.lean}). A commutative algebra has a distributive
(Boolean) event lattice by the Birkhoff--von~Neumann correspondence, so a joint
probability distribution over its observables \emph{always} exists. In a word:
\textbf{the substrate's logic is classical logic.} Non-contextual hidden
variables are not an extra hypothesis here---they are forced by the algebra. The
full argument, and the crucial Poisson-versus-commutator distinction it turns
on, is reproduced in \S\ref{sec:derivation}.

\section{The four closed routes to non-commutativity}\label{sec:closed-routes}

Result~\ref{res:commutative} is an \emph{absence} statement, and absence
statements earn their keep only by surviving attempts to circumvent them. Each
known FTD-native route to manufacturing quantum non-commutativity from P1--P5
has been tested and closed:

\begingroup
\renewcommand{\arraystretch}{1.3}
\begin{center}
\begin{tabular}{@{}>{\raggedright}p{0.38\textwidth} >{\raggedright}p{0.36\textwidth} >{\raggedright\arraybackslash}p{0.15\textwidth}@{}}
\toprule
\textbf{Route} & \textbf{Why it closes} & \textbf{Status} \\
\midrule
Modular time --- a Tomita--Takesaki type~III$_1$ algebra read off the substrate
  & The substrate algebra is commutative $\to$ type~I $\to$ trivial modular flow
  & \ctag{cn}\newline\ftdid{0225} \\
Manifestation as measurement --- genesis $+$ Gauss back-reaction read as
  collapse
  & A deterministic function of commuting flux gives a Boolean event lattice,
    i.e.\ classical coarse-graining
  & \ctag{cn}\newline\ftdid{0226} \\
Budget symmetry --- the body-diagonal $C_3$ (the $N_c=3$ count) read as the
  epistemic $\mathbb{Z}/3$
  & $\{J_x,J_y,J_z\}$ \emph{commute}: a co-measurable triple, not three
    complementary bases --- same count, wrong kind
  & \ctag{cn}\newline\ftdid{0228} \\
The Bell wall --- a substrate-native CHSH violation $>2$
  & A local deterministic substrate obeys $S\le 2$ (Bell's theorem applies);
    $S=2\sqrt2$ appears only \emph{after} the complexification
    $J_x+iJ_y\mapsto\psi$, which \emph{is} an instance of $M$
  & \ctag{cn}\newline\ftdid{0243}~\S4 \\
\bottomrule
\end{tabular}
\end{center}
\endgroup

\noindent
The four closures share a moral: where a candidate for non-commutativity seems
to appear, it is always either a co-measurable structure wearing the wrong name,
or it is the externally imported measurement map $M$ smuggled in by a
complexification step. The substrate itself stays commutative throughout.

\section{FC-1 --- the framework declines the measurement map $M$}\label{sec:fc1}

\ftdid{0243} establishes more than the commutativity result: it establishes that
non-commutativity is \emph{logically independent} of P1--P5. Both
$\{\mathrm{P1\ldots P5}\}\cup\{M\}$ and $\{\mathrm{P1\ldots P5}\}\cup\{\neg M\}$
are consistent, where $M$ is an independent measurement / quantization postulate
(a complexification and projection-valued measure, or a 't~Hooft template-state
basis). The theorem proves the \emph{fork exists}. A theorem cannot choose which
branch of a genuine fork is true; only a commitment can.

\begin{conditionalbox}[FC-1 (framework commitment)]
The commutative observable algebra $A_5$ is \emph{complete}: FTD declines the
measurement-map import $M$. Non-commutativity, Hilbert-space state vectors, and
operator-valued observables are not part of FTD's model of the world. Where
quantum mechanics' non-commutative formalism makes structural predictions that
differ from the commutative substrate's, FTD predicts the substrate.
\end{conditionalbox}

\noindent
\etag{ax}~(\ftdid{0255}, an \etag{ax}-class declaration; licensed but not entailed
by \ftdid{0243}). \textbf{This is a declaration, not a derivation.} The theorem
proves the fork; FC-1 picks the $\neg M$ branch. Three consequences follow, each
taken up later in the chain: measurement is re-located to the observer layer as a
frame-relative epistemic restriction (Chapter~\ref{ch:philosophy}); the wave-face
of quantum phenomenology is carried by the classical flux (Chapter~\ref{ch:physics});
and the structural \emph{differences} between the commutative substrate and the
quantum formalism become FTD's falsifiable spine
(Chapter~\ref{ch:science}). FC-1 places FTD in the $\psi$-epistemic tradition (a
classical ontic substrate with epistemic restrictions) while going beyond it:
epistricted theories aim to \emph{reproduce} a fragment of QM, whereas FTD
declines the reproduction target itself. What would overturn FC-1 is stated in
\S\ref{sec:boundary}.

\section{FC-2 --- the arrow of time is native}\label{sec:fc2}

The same logical base supports a second declaration. \ftdid{0253} maps the gap:
reversibility is absent from P1--P5, and the genuinely \emph{finite} sector of
the substrate---the ternary state field---is irreversible, since manifestation
is many-to-one and provably unrecoverable. The reversible part of the dynamics
is the continuous flux, on which the finiteness commitment has no purchase.
Finiteness therefore actively \emph{opposes} a global reversibility postulate
rather than merely failing to imply one.

\begin{conditionalbox}[FC-2 (framework commitment, arrow half)]
The arrow of time is native: the manifestation/state sector is fundamentally
irreversible, and FTD declines global reversibility as a postulate.
Reversibility holds only as a \emph{sector} property of the weakly-coupled flux
wave dynamics---and with it everything that rides on reversibility, including the
Lorentzian metric.
\end{conditionalbox}

\noindent
\etag{ax}~(\ftdid{0256}, an \etag{ax}-class declaration; licensed by \ftdid{0253}
and the scoped IR measurement \ftdid{0252}). Again \textbf{a declaration, not a
derivation}: \ftdid{0253} proves only that the postulates do not \emph{force}
reversibility, and that FTD's finiteness argues against it; FC-2 commits to the
irreversible reading as the framework's official model. Two scope notes keep this
honest. First, the causal cone $\cwave$ remains forced (\etag{thm}, from P4); it
is the metric structure layered on top of the cone---and the Lorentzian
space--time mixing it carries---that is declared emergent and sector-scoped, not
the cone itself. Second, space (P1) and time (P2) are fundamentally separate
primitives at the base: substrate simultaneity is absolute, and Minkowski mixing
is an infrared property of the flux wave sector only. The metric half of FC-2 is
developed where it belongs, in Chapter~\ref{ch:physics}; here Logic records only
the arrow half and its declaration status. What would overturn FC-2 is stated in
\S\ref{sec:boundary}.

\section{The epistemic horizon, honestly split}\label{sec:horizon}

If the substrate's logic is classical, why does observing it feel
restricted---why can an internal observer not simply read the whole
configuration? FTD's answer splits cleanly into a part it can prove and a part
it cannot, and the honesty of the split matters more than either half.

\begin{description}[leftmargin=1.4em,itemsep=4pt]
\item[The binding half.] A finite internal observer---a subsystem with $M$
  pointer states attempting to model a total of $N>M$ states, \emph{itself
  included}---necessarily loses access to part of its own state. This is
  ordinary classical finite self-reference (a pigeonhole argument); it uses no
  quantum input and is non-circular. Every internal frame therefore carries a
  \textbf{self-blind-spot}. \etag[ binding half]{thm}~(\ftdid{0227}).
\item[The sharp half.] The \emph{Spekkens knowledge-balance}---the blind-spot
  rotating symmetrically over complementary bases, which is what would yield
  genuinely quantum-like epistemic statistics---requires the full symplectic
  $S_3 = \mathbb{Z}/2 \rtimes \mathbb{Z}/3$ action. The $\mathbb{Z}/2$ is
  FTD-native (the $J^2=-I$ quarter-conjugacy); the $\mathbb{Z}/3$ is
  \emph{missing} (the $N_c=3$ candidate was the same co-measurable triple closed
  as route~3 above). Binding-without-sharpness is the framework's honest state.
  \etag{opn}~(\ftdid{0227}).
\end{description}

\noindent
A post-hoc note, flagged as such: binding-without-sharpness \emph{after the
fact} accounts for the engine's measured detection statistics being Rice /
Gaussian rather than Born (the PL-1 row of Chapter~\ref{ch:science})---a classical
restriction without the symmetric balance gives classical-with-noise statistics,
not $|\psi|^2$. This account was constructed after the measurement; it is an
explanation candidate, not a pre-registered prediction, and it is tagged
accordingly.

% ---- the reproduced derivation ---------------------------------------------
\section{Reproduced derivation: commutativity independence}\label{sec:derivation}

This is the load-bearing argument of the layer. It is reproduced in full because
the entire weight of FC-1 rests on it, and because its crucial step---that a
non-trivial symplectic structure and a commutative observable algebra coexist
without contradiction---is exactly the point an outside reader is most likely to
doubt. The argument has two claims (absence, then independence), with the
Poisson-versus-commutator distinction sitting between them.

\subsection*{Claim A --- absence: $A_5$ is strictly commutative}

Every generator in the inventory $\{s, J_a, \mathrm{wave\_vel}, \mathcal{L}\}$
is a real-valued function on the single configuration space $\Omega$ (an
assignment of $(s(v), J(v), \mathrm{wave\_vel}(v), \mathcal{L}(v))$ to every
voxel $v$). The algebra's multiplication is the \emph{pointwise} product, and
real multiplication is trivially commutative: for all $\omega\in\Omega$,
\[
  (A\cdot B)(\omega) \;=\; A(\omega)\,B(\omega) \;=\; B(\omega)\,A(\omega)
  \;=\; (B\cdot A)(\omega).
\]
The closure operations preserve function-hood: a Moore-neighbourhood sum of
functions is a function, and a composition $A\circ U$ with the deterministic
update map is again a function $\Omega\to\mathbb{R}$. Hence every element of
$A_5$ is a real-valued function on $\Omega$, and the product of any two commutes
pointwise. Therefore $A_5$ is a commutative real algebra. (The algebraic core is
machine-checked: \texttt{observable\_commutator\_zero} holds for all pointwise
configurations in \texttt{lean/Standalone.lean}.) \hfill$\square$

\subsection*{The crucial distinction: $\{q,p\}\neq 0$ yet $[q,p]=0$}

Here is the step that defuses the natural objection. The leapfrog / symplectic
time-update of the substrate endows its phase space with a \emph{non-zero
Poisson bracket},
\[
  \{q,p\} \;=\; 1,
\]
where $(q,p)=(J,\mathrm{wave\_vel})$ are the substrate's quadratures. The Poisson
bracket is an antisymmetric bilinear \emph{derivation on phase-space functions}.
It is a real structure on the \emph{space of beables}---it is what makes the
quadrature phase wind, and it is genuinely non-trivial. It is \emph{not} the
algebra's multiplication.

The observable commutator is a different object entirely. It is built from the
\emph{pointwise product}, and by Claim~A that product is symmetric, so
\[
  [q,p] \;=\; q\cdot p - p\cdot q \;=\; 0
\]
identically, on every configuration. The two facts
\begin{equation}\label{eq:crucial}
  \boxed{\;\{q,p\}\neq 0 \quad\text{yet}\quad [q,p]=0\;}
\end{equation}
are therefore perfectly consistent: $\{\cdot,\cdot\}$ acts on functions,
$[\cdot,\cdot]$ acts on the algebra's product, and a structure on the first does
not imply a structure on the second. \etag{thm}~(\ftdid{0243}).

The intuition worth carrying away: the substrate is \emph{symplectically rich}
and \emph{algebraically commutative} simultaneously. Its quadratures wind---there
is a genuine clock---but all quadratures \emph{co-measure}; choosing one as
``the'' measured value, and declaring the conjugate one disturbed, is not a fact
the substrate provides. What would convert the non-trivial Poisson structure into
a non-commutative operator algebra is a \emph{deformation quantization}, a map
$\{\cdot,\cdot\}\mapsto[\cdot,\cdot]/i\hbar$. That map intrinsically requires the
external deformation parameter $\hbar$---it acts as precisely the external
measurement map $M$---and is not derivable from $A_5$ alone. The symplectic
structure invites quantization; it does not perform it.

\subsection*{Claim B --- independence: $M$ cannot be distilled from $A_5$}

By Claim~A, $A_5$ is strictly commutative, so its event lattice is distributive
(Boolean) and a joint probability distribution always exists. Any standard QFT or
relativistic-time target strictly requires a non-commuting pair, $[A,B]\neq 0$
(a non-distributive event lattice / type~III$_1$ modular algebra).
Commutativity, however, is algebraically \emph{preserved} under topological
limits and classical coarse-graining: a limit of commuting families commutes, and
coarse-graining a Boolean lattice yields a Boolean lattice. Therefore no limit
and no coarse-graining applied to $A_5$ can ever organically yield
$[A,B]\neq 0$. The non-commutative structure is logically independent of P1--P5
and must be supplied by an external map $M$. That
$\{\mathrm{P1\ldots P5}\}\cup\{M\}$ is consistent is witnessed by the existing
emergent-QM constructions that explicitly introduce such an $M$; consistency
together with non-derivability is exactly logical independence. \hfill$\square$

\subsection*{Why the Bell violation does not contradict this}

A consistency check, because the framework does exhibit a Bell violation. The
$S=2\sqrt2$ value in FTD is \emph{emergent}, derived via the Gauss constraint as
a transverse non-local bounding mechanism: the CHSH-violating correlations arise
\emph{geometrically} from that constraint, not from native non-commuting
measurement operators. But translating those emergent classical correlations into
the standard QM Bell format relies on a complexification
$J_x + iJ_y \mapsto \psi$ with coarse-graining---and that complexification
\emph{is} an instance of the independent map $M$. The substrate itself remains
strictly commutative and limited to $S\le 2$; the apparent $2\sqrt2$ emerges only
once $M$ is applied. The theorem holds, sharpening the boundary without
contradicting the emergent result. (The unresolved laboratory burden this leaves
is PL-2 in Chapter~\ref{ch:science}, the framework's highest-risk row.)

\subsection*{What the argument licenses}

Claims A and B together are the licensing theorem for FC-1: they prove the fork
is genuine (both branches consistent) and that the substrate cannot cross it on
its own. They do \emph{not} pick a branch---that is FC-1's job, and FC-1 is a
declaration. The same argument supplies the parallel that organizes the rest of
the chain: the relativistic \emph{metric} needs reversibility (declined by FC-2)
and the quantum \emph{structure} needs non-commutativity (declined by FC-1), and
these are \emph{the same shape of gap}---each a single added principle on the
same classical, commutative, irreversible base.

% ---- boundary subsection ----------------------------------------------------
\section{What this layer does not settle}\label{sec:boundary}

Logic's two strongest moves are \emph{declarations}, and the discipline of the
chain is to say so and to state in advance what would overturn each. None of the
following is softened; the kill-conditions are quoted from the constitution's
falsification register.

\paragraph{FC-1 and FC-2 are declarations, not theorems.} The licensing results
(\ftdid{0243} for FC-1, \ftdid{0253} for FC-2) prove only that the postulates
leave the relevant fork open. Picking the branch is a model-defining choice,
tagged \etag{ax}-class and never to be read as derived. A reader who accepts the
postulates is \emph{not} thereby compelled to accept either commitment; that is
the honest cost of taking them.

\paragraph{What kills FC-1.} Any one of: (i) a substrate-native pair of
observables with $[A,B]\neq 0$ derived from P1--P5 (this would refute
\ftdid{0243}'s premise and is the single most direct refutation of the
framework); (ii) a P1--P5 derivation of Born statistics; (iii) a substrate-native,
$M$-free correlation experiment exceeding $S=2$ beyond numerical error.

\paragraph{What kills FC-2.} Any one of: (i) exact $\gamma$ (exact Lorentz
invariance) at finite $L$ off the infrared limit---a measured \emph{absence} of
the predicted ultraviolet corrections where the lattice requires them; (ii) a
P1--P5 derivation of global reversibility (this would refute \ftdid{0253}'s
boundary); (iii) failure of the measured IR approach---the $L^{-2}$ law breaking
on the $\langle 100\rangle$ axis at larger $L$, or the $k^4$ anisotropy decay
reversing. The detailed kill-conditions and their protocols are the subject of
Chapter~\ref{ch:science}.

\paragraph{The sharpness of the blind spot is open.} The binding half of the
epistemic horizon is a theorem; the \emph{sharp} half---the symmetric
knowledge-balance that would make the observer's restriction genuinely
quantum-like---is unresolved, blocked on the missing $\mathbb{Z}/3$ of the
symplectic budget symmetry. \etag{opn}~(\ftdid{0227}). Binding-without-sharpness
is where the observer story honestly stands; it is carried forward, not papered
over, into the philosophy layer (Chapter~\ref{ch:philosophy}).
